\documentclass[10pt]{article}
\usepackage{cornell-notes}
\title{Chapter 12: Logic Interfacing}\author{}\date{}
\setDocTitle{Chapter 12: Logic Interfacing}\setDocSubtitle{Electronics}\setDocOwner{Electronics Cornell Notes}\setDocDate{\today}
\setCornellCollection{Electronics Cornell Notes}\setCornellUnitType{Chapter}\setCornellUnitNumber{12}\setCornellUnitTitle{Logic Interfacing}
\hypersetup{pdftitle={Chapter 12: Logic Interfacing},pdfsubject={Cornell notes},pdfauthor={}}
\begin{document}\maketitle
\begin{cornell}\noterow{Core idea}{Interfaces reconcile voltage domains, timing, capacitance, protection, and signal integrity.}\noterow{Validation}{Check thresholds, drive strength, slew rate, termination, and fault conditions.}\end{cornell}
\begin{summarybox}{Section summary}An interface must be evaluated at both sides of the electrical boundary.\end{summarybox}
\end{document}
